\documentclass[10pt]{extarticle}
\usepackage[a4paper,margin=1.5cm]{geometry}
\usepackage[T1]{fontenc}
\usepackage{amsmath,amssymb}
\usepackage{graphicx}
\usepackage[export]{adjustbox}
\usepackage{multicol}
\usepackage{enumitem}
\usepackage{booktabs,array,tabularx}
\setlength{\columnsep}{0.7cm}
\setlength{\parindent}{0pt}
\setlength{\parskip}{3pt}
\pagestyle{empty}
\makeatletter
% figure registry (figdefs.tex is generated by the extract step)
\newcommand{\qfigdef}[4]{\@namedef{qf@file@#1}{#2}\@namedef{qf@w@#1}{#3}\@namedef{qf@wide@#1}{#4}}
\newcommand{\qf@check}[1]{\@ifundefined{qf@file@#1}{\PackageError{score2}{Unknown figure id #1}{}}{}}
\newcommand{\qf@img}[4][]{\includegraphics[width=#2\dimexpr\@nameuse{qf@w@#4}pt\relax,max width=#3,max height=0.45\textheight,#1]{\@nameuse{qf@file@#4}}}
% Full-width questions: finalize wraps every question that contains a wide figure in \qfullbegin...\qfullend,
% so the WHOLE question (stem, figures, options) is set across the page in reading order, then the two columns
% resume. The question's head - its number, stem and first figure (\qfullsplit marks the end of the head) -
% is kept on one page, so a number is never left at a page foot with its figure overleaf; the rest may break
% across pages normally, so no big gaps are left. Without a \qfullsplit, the whole question is kept together
% if it is under 45% of a page. Never write these yourself.
\newif\ifqf@full
\newif\ifqf@boxed
\newsavebox{\qf@box}
\newcommand{\qf@open}{\global\qf@fulltrue\par\global\qf@boxedtrue
  \setbox\qf@box\vbox\bgroup\hsize\textwidth\linewidth\textwidth}
\newcommand{\qf@emit}[1]{\par\egroup\global\qf@boxedfalse
  \ifdim\dimexpr\ht\qf@box+\dp\qf@box\relax<#1\textheight\noindent\usebox{\qf@box}\par\else\unvbox\qf@box\fi}
\newcommand{\qfullsplit}{\ifqf@boxed\qf@emit{0.85}\fi}
\newcommand{\qf@close}{\ifqf@boxed\qf@emit{0.45}\fi\global\qf@fullfalse}
\newcommand{\qfullbegin}{\end{multicols}\qf@open}
\newcommand{\qfullend}{\qf@close\begin{multicols}{2}}
\newcommand{\qfullstart}{\qf@open}
\newcommand{\qfulldone}{\qf@close}
% \qfig[scale]{id}: block figure, centred, directly where it is written (size set by the tool).
% Fallback if a wide figure is not inside a full-width question: break out of the columns in place - no floats.
\newcommand{\qfig}[2][1]{\qf@check{#2}%
  \ifqf@full\qf@blk{#1}{#2}\else
  \ifnum\@nameuse{qf@wide@#2}=1 \expandafter\@firstoftwo\else\expandafter\@secondoftwo\fi
  {\end{multicols}\begin{center}\qf@img{#1}{\textwidth}{#2}\end{center}\begin{multicols}{2}}%
  {\qf@blk{#1}{#2}}\fi}
\newcommand{\qf@blk}[2]{\par\smallskip{\centering\qf@img{#1}{\linewidth}{#2}\par}\smallskip}
% \qfigi[scale]{id}: inline figure (inside option lists / tables / minipages); never breaks columns.
\newcommand{\qfigi}[2][1]{\qf@check{#2}\qf@img[valign=c]{#1}{\linewidth}{#2}}
\makeatother
% structure macros - lint reads these to check completeness and figure placement
\newcommand{\papersec}[1]{\par\bigskip{\centering\Large\bfseries PAPER-#1\par}\medskip}
\newcommand{\partsec}[2]{\par\medskip\noindent{\large\bfseries PART-#1: #2}\par\smallskip}
\newcommand{\sect}[1]{\par\smallskip\noindent\textbf{#1}\par}
\newcommand{\qn}[1]{\par\smallskip\noindent\textbf{#1.}~}
\newcommand{\paraq}[2]{\par\smallskip\noindent\textbf{Paragraph for Questions #1 and #2}\par}
\newenvironment{qoptions}{\begin{itemize}[leftmargin=1.8em,itemsep=1pt,topsep=2pt,label={}]}{\end{itemize}}
\newcommand{\opt}[2]{\item[(#1)] #2}
% figure registry (generated by score2_tools.py finalize)
\qfigdef{p1_phy_s1_q01_1}{TEAS - Test 01 (16 Apr 2023)-img/p1_phy_s1_q01_1.png}{161.3}{0}
\qfigdef{p1_phy_s1_q02_1}{TEAS - Test 01 (16 Apr 2023)-img/p1_phy_s1_q02_1.png}{153.3}{0}
\qfigdef{p1_phy_s1_q03_1}{TEAS - Test 01 (16 Apr 2023)-img/p1_phy_s1_q03_1.png}{148.0}{0}
\qfigdef{p1_phy_s1_q04_1}{TEAS - Test 01 (16 Apr 2023)-img/p1_phy_s1_q04_1.png}{161.3}{0}
\qfigdef{p1_phy_s1_q05_1}{TEAS - Test 01 (16 Apr 2023)-img/p1_phy_s1_q05_1.png}{107.8}{0}
\qfigdef{p1_phy_s2_para7_8_1}{TEAS - Test 01 (16 Apr 2023)-img/p1_phy_s2_para7_8_1.png}{115.1}{0}
\qfigdef{p1_phy_s2_para9_10_1}{TEAS - Test 01 (16 Apr 2023)-img/p1_phy_s2_para9_10_1.png}{66.3}{0}
\qfigdef{p1_phy_s3_q01_1}{TEAS - Test 01 (16 Apr 2023)-img/p1_phy_s3_q01_1.png}{160.9}{0}
\qfigdef{p1_phy_s3_q02_1}{TEAS - Test 01 (16 Apr 2023)-img/p1_phy_s3_q02_1.png}{108.3}{0}
\qfigdef{p1_phy_s3_q04_1}{TEAS - Test 01 (16 Apr 2023)-img/p1_phy_s3_q04_1.png}{92.6}{0}
\qfigdef{p1_phy_s3_q05_1}{TEAS - Test 01 (16 Apr 2023)-img/p1_phy_s3_q05_1.png}{107.8}{0}
\qfigdef{p1_phy_s3_q06_1}{TEAS - Test 01 (16 Apr 2023)-img/p1_phy_s3_q06_1.png}{87.4}{0}
\qfigdef{p1_phy_s3_q08_1}{TEAS - Test 01 (16 Apr 2023)-img/p1_phy_s3_q08_1.png}{213.9}{0}
\qfigdef{p1_chem_s1_q01_1}{TEAS - Test 01 (16 Apr 2023)-img/p1_chem_s1_q01_1.png}{196.2}{0}
\qfigdef{p1_chem_s1_q04_A}{TEAS - Test 01 (16 Apr 2023)-img/p1_chem_s1_q04_A.png}{92.5}{0}
\qfigdef{p1_chem_s1_q04_B}{TEAS - Test 01 (16 Apr 2023)-img/p1_chem_s1_q04_B.png}{78.2}{0}
\qfigdef{p1_chem_s1_q04_C}{TEAS - Test 01 (16 Apr 2023)-img/p1_chem_s1_q04_C.png}{90.9}{0}
\qfigdef{p1_chem_s1_q04_D}{TEAS - Test 01 (16 Apr 2023)-img/p1_chem_s1_q04_D.png}{34.2}{0}
\qfigdef{p1_chem_s2_para7_8_1}{TEAS - Test 01 (16 Apr 2023)-img/p1_chem_s2_para7_8_1.png}{238.3}{0}
\qfigdef{p1_chem_s2_para7_8_2}{TEAS - Test 01 (16 Apr 2023)-img/p1_chem_s2_para7_8_2.png}{206.9}{0}
\qfigdef{p1_chem_s2_para7_8_3}{TEAS - Test 01 (16 Apr 2023)-img/p1_chem_s2_para7_8_3.png}{273.1}{1}
\qfigdef{p1_chem_s2_para7_8_4}{TEAS - Test 01 (16 Apr 2023)-img/p1_chem_s2_para7_8_4.png}{288.9}{1}
\qfigdef{p1_chem_s2_q07_1}{TEAS - Test 01 (16 Apr 2023)-img/p1_chem_s2_q07_1.png}{138.1}{0}
\qfigdef{p1_chem_s2_q07_A}{TEAS - Test 01 (16 Apr 2023)-img/p1_chem_s2_q07_A.png}{78.7}{0}
\qfigdef{p1_chem_s2_q07_B}{TEAS - Test 01 (16 Apr 2023)-img/p1_chem_s2_q07_B.png}{91.1}{0}
\qfigdef{p1_chem_s2_q07_C}{TEAS - Test 01 (16 Apr 2023)-img/p1_chem_s2_q07_C.png}{79.9}{0}
\qfigdef{p1_chem_s2_q07_D}{TEAS - Test 01 (16 Apr 2023)-img/p1_chem_s2_q07_D.png}{95.4}{0}
\qfigdef{p1_chem_s2_q08_1}{TEAS - Test 01 (16 Apr 2023)-img/p1_chem_s2_q08_1.png}{244.4}{0}
\qfigdef{p1_chem_s3_q06_1}{TEAS - Test 01 (16 Apr 2023)-img/p1_chem_s3_q06_1.png}{74.4}{0}
\qfigdef{p1_chem_s3_q06_2}{TEAS - Test 01 (16 Apr 2023)-img/p1_chem_s3_q06_2.png}{107.5}{0}
\qfigdef{p1_chem_s3_q07_1}{TEAS - Test 01 (16 Apr 2023)-img/p1_chem_s3_q07_1.png}{191.3}{0}
\raggedcolumns
\makeatletter\@beginparpenalty=10000\makeatother
\newcolumntype{L}[1]{>{\raggedright\arraybackslash}p{#1}}
\begin{document}
\begin{center}{\Large\bfseries TEAS --- Test 01 (16 Apr 2023)}\end{center}
\begin{multicols}{2}
\raggedcolumns
\partsec{1}{PHYSICS}
\sect{SECTION-I(i)}
\qn{1} Magnetic field in a cylindrical region of radius `$R$' is given as $B=B_0t$. Direction of magnetic field is out of plane. A point charge `$q$' is slowly moved from origin to infinity. Choose the correct statement regarding work done by external agent.
\qfig{p1_phy_s1_q01_1}
\begin{qoptions}
\opt{A}{Work done is zero, if it is moved along $x$-axis.}
\opt{B}{Work done is $\dfrac{\pi R^2}4B_0q$, if it is moved along $y$-axis.}
\opt{C}{Work done is independent of direction of movement of charge particle.}
\opt{D}{Work done is $\dfrac{\pi R^2B_0q}8$, if it is moved along $y$-axis.}
\end{qoptions}
\qn{2} A calorimeter was filled with 10 g of water at a temperature of 70$^\circ$C. To cool it down, the researchers threw in calorimeter some crushed ice from a freezer at temperature $-10\,^\circ$C. Latent heat of ice $=330$ kJ/kg, specific heat capacity of water and ice are $c_w=4200$ J/(kg $^\circ$C) and $c_i=2100$ J/(kg $^\circ$C) respectively Graphs of dependences of the mass of each of substances from time are given on the image. Unfortunately, scientists have spilled coffee on the graph, and part of the experimental data was lost. Help researchers recover these data by choosing the correct answers. Neglect any heat loss.
\qfig{p1_phy_s1_q02_1}
\begin{qoptions}
\opt{A}{The temperature of the water by the time the ice began to melt was 60$^\circ$C.}
\opt{B}{The final temperature in the calorimeter was 0$^\circ$C.}
\opt{C}{In steady state, 50\% was water and rest ice.}
\opt{D}{The mass of ice initially was 20 gm.}
\end{qoptions}
\qn{3} The figure shows a diagram of an electrical circuit. All resistors are the same resistance $R=1$ Ohm. Connect between points $A$ and $B$ ideal battery with a voltage of 8 V. Find the readings of the ideal ammeters.
\qfig{p1_phy_s1_q03_1}
\begin{qoptions}
\opt{A}{$A_1=3$A}
\opt{B}{$A_2=0$A}
\opt{C}{$A_1=1$A}
\opt{D}{$A_2=1$A}
\end{qoptions}
\qn{4} In the diagram we have three capacitors $C_1=20\ \mu$F, $C_2=30\ \mu$F and $C_3=60\ \mu$F, two sources with emf $\varepsilon_1=1$ V and $\varepsilon_2=2$ V. The potential difference $V_A-V_B=3$ V. Find the voltage on each capacitor.
\qfig{p1_phy_s1_q04_1}
\begin{qoptions}
\opt{A}{$V_1=1$ volt}
\opt{B}{$V_2=2/3$ volt}
\opt{C}{$V_3=1/3$ volt}
\opt{D}{$V_1=2$ volt}
\end{qoptions}
\qn{5} A point charge $q$, mass $m$ is projected at an angle $\theta$ to a level ground in presence of uniform magnetic field $B$. Gravity is absent. The collision of particle with ground is elastic choose the correct statement(s). ($B=0.05$ T, $q=1\ \mu$C, $m=4\times10^{-5}$ kg, $\theta=30^\circ$, $v=1$ cm/s)
\qfig{p1_phy_s1_q05_1}
\begin{qoptions}
\opt{A}{The time duration between 2 collisions is $\dfrac{800\pi}3$ sec.}
\opt{B}{The distance between 2 successive collision is 8m.}
\opt{C}{Maximum height from ground is 4m.}
\opt{D}{Average force on ground after a long time is $\dfrac{3\pi}2\times10^{-9}$N.}
\end{qoptions}
\qn{6} The radiation coming from a black body at a temperature of 3000 K are analyzed. Assume wein's constant $b=3\times10^{-3}$ mK. We wish to compare intensity of radiation corresponding to wavelength range (999 nm to 1001 nm), say $E_1$; wavelength range (1499 nm to 1501 nm), say $E_2$; wavelength range (499 nm to 501 nm) say $E_3$.
\begin{qoptions}
\opt{A}{$E_2>E_1$}
\opt{B}{$E_2>E_3$}
\opt{C}{$E_1>E_3$}
\opt{D}{$E_3>E_1$}
\end{qoptions}
\sect{SECTION-I(ii)}
\paraq{7}{8}
In the given circuit, capacitance is $\dfrac{500}\pi\ \mu$F and inductance is $\dfrac1{5\pi}$ H.
\qfig{p1_phy_s2_para7_8_1}
\qn{7} Current $I_1$ in the given circuit is :-
\begin{qoptions}
\opt{A}{$\sqrt5$A}
\opt{B}{$\sqrt{\dfrac52}$A}
\opt{C}{$\sqrt2$A}
\opt{D}{$\dfrac1{\sqrt2}$A}
\end{qoptions}
\qn{8} Power factor of the given circuit is :-
\begin{qoptions}
\opt{A}{$\dfrac2{\sqrt5}$}
\opt{B}{$\dfrac1{\sqrt2}$}
\opt{C}{$\dfrac12$}
\opt{D}{1}
\end{qoptions}
\paraq{9}{10}
A non-conducting thin disc has surface charge density $\sigma=\dfrac{\sigma_0r}R$. Disc has radius `$R$' and being rotated with angular velocity $\omega$.
\qfig{p1_phy_s2_para9_10_1}
\qn{9} Magnetic moment of disc is :-
\begin{qoptions}
\opt{A}{$\dfrac{\sigma_0\omega\pi R^4}4$}
\opt{B}{$\dfrac{\sigma_0\omega R^4}4$}
\opt{C}{$\dfrac{\sigma_0\omega\pi R^4}5$}
\opt{D}{$\dfrac{\sigma_0\omega\pi R^4}2$}
\end{qoptions}
\qn{10} Magnetic field at the centre of disc is :-
\begin{qoptions}
\opt{A}{$\dfrac{\mu_0\sigma_0\omega R}4$}
\opt{B}{$\dfrac{\mu_0\sigma_0\omega R}2$}
\opt{C}{$\mu_0\sigma_0\omega R$}
\opt{D}{$4\mu_0\sigma_0\omega R$}
\end{qoptions}
\sect{SECTION-II}
\qn{1} Find the average power dissipated in the circuit in kilo-watt.
\qfig{p1_phy_s3_q01_1}
\qn{2} In the given circuit, current through battery is independent of time. If the resistance `$R$' to satisfy this condition is $25K$ then find `$K$'.
\qfig{p1_phy_s3_q02_1}
\qn{3} A conducting ring has radius `$R$'. A magnetic dipole is moved along the axis of ring such that total flux changed through ring is $\phi$. Area of cross section of ring is `$A$' and number density of conducting electron is `$n$'. Current flowing through ring is given by $I=\dfrac{ne^2\phi A}{K\pi mR}$. Find `$K$'. (The time taken for this flux change is equal to relaxation time.) [Here $e=$ charge of electron, $m=$ mass of electron]
\qn{4} Number of turns per unit length is `$n$' and current through solenoid is `$i$'. Radius of the solenoid is $\dfrac\ell3$. If magnetic field on the axis of solenoid at point `$P$' is $\dfrac{\mu_0ni}{2\sqrt{10}}\left[\sqrt y+x\sqrt x\right]$ the find $x+y$.
\qfig{p1_phy_s3_q04_1}
\qn{5} In the given figure resistance of potentiometer wire is 10$\Omega$. If the balancing length is $\ell_1$ and $\ell_2$ when switch is open and closed then find $\dfrac{\ell_1}{\ell_2}$.
\qfig{p1_phy_s3_q05_1}
\qn{6} Magnetic field in a cylindrical region centred at origin is given as $\vec B=\dfrac13t$ Tesla $(-\hat k)$. Find the emf induced in the the rod $AB$ kept in $x$-$y$ plane as shown in figure.
\qfig{p1_phy_s3_q06_1}
\qn{7} In a thermodynamic prcess, temperature of gas as function of entropy is given as $T=\alpha.S^2$ where `$T$' is temperature and `$S$' is entropy of gas `$\alpha$' is a positive constant. If heat capacity of gas as a function of entrpy is $C=\dfrac Sx$ then find `$x$'.
\qn{8} Two rings of equal mass `$m$' are separated by `$r$' distance initially and released from there. They carry same current `$i$' as shown. If velocity of both the rings when they are very far away from each other is $\dfrac{ir_1r_2k}{2r}\sqrt{\dfrac{\mu_0\pi}{mr}}$. Find `$k$'. Assume that $r_1$ and $r_2$ are much smaller than $r$.
\qfig{p1_phy_s3_q08_1}
\partsec{2}{CHEMISTRY}
\sect{SECTION-I(i)}
\qn{1} Consider the following reaction
\qfig{p1_chem_s1_q01_1}
(I) $(\mathrm{IE})_D<(\mathrm{IE})_{\text{oxygen atom}}$\\
(II) $E+\mathrm{BF_3}\to$ Monobasic acid\\
(III) $A$ is non-polar \& planar\\
(IV) $C$ contains $\mathrm{p_\pi}$-$\mathrm{d_\pi}$ bonds\\
Choose the correct statement(s) on the basis of above information
\begin{qoptions}
\opt{A}{Number of electrons in gerade molecular orbitals in ``$D$'' is 8}
\opt{B}{``$B$'' can be liquified with more ease compared to $\mathrm{NH_3}$}
\opt{C}{``$C$'' is polar \& non planar}
\opt{D}{``$C$'' has distorted octahedral shape}
\end{qoptions}
\qn{2} Which of the following inner orbital complex(es) is/are paramagnetic.
\begin{qoptions}
\opt{A}{$\mathrm{Na_3[CoF_6]}$}
\opt{B}{$\mathrm{[Cr(H_2O)_6]^{3+}}$}
\opt{C}{$\mathrm{K[PtCl_3(\pi\text{-}C_2H_4)]}$}
\opt{D}{$\mathrm{[Co(en)_3]Cl_2}$}
\end{qoptions}
\qn{3} Which of the following statement(s) is/are correct.
\begin{qoptions}
\opt{A}{According to Werner theory, the ions/groups bounded by primary valency to the metal have characteristic spatial arrangements corresponding to different coordination number}
\opt{B}{Spectrochemical series is based on absorption of light by complexes for a particular metal ion with different ligands}
\opt{C}{$\mathrm{d^4}$ to $\mathrm{d^7}$ coordination entities are more stable for weak field ligands as compared to strong field ligands}
\opt{D}{The average energy of a metal ion in a spherical field is higher as compared to average energy of isolated free metal ion.}
\end{qoptions}
\qn{4} Which of the following reaction result the formation of a pair of diastereomers?
\begin{qoptions}
\opt{A}{\qfigi{p1_chem_s1_q04_A}}
\opt{B}{\qfigi{p1_chem_s1_q04_B}}
\opt{C}{\qfigi{p1_chem_s1_q04_C}}
\opt{D}{\qfigi{p1_chem_s1_q04_D}}
\end{qoptions}
\qn{5} 8.0 gm of solid sulphur is first oxidised to $\mathrm{SO_2}$ and then it is divided into two equal parts. One part is sufficient for complete decolorisation of 200 ml of acidified $\mathrm{KMnO_4}$ solution. Another part is oxidised to $\mathrm{SO_3}$ and the $\mathrm{SO_3}$ formed is sufficient for complete precipitation of all $\mathrm{BaCl_2}$ present in 100 ml solution as $\mathrm{BaSO_4}$. Which of the following statements is/are correct ? ($\mathrm S=32$, $\mathrm{Ba}=138$)
\begin{qoptions}
\opt{A}{The molarity of $\mathrm{KMnO_4}$ solution is 0.25}
\opt{B}{The molarity of $\mathrm{BaCl_2}$ solution is 0.25}
\opt{C}{The weight of $\mathrm{BaSO_4}$ precipitated out is 29.25 gm.}
\opt{D}{The same moles of $\mathrm{KMnO_4}$ and $\mathrm{BaCl_2}$ are reacted.}
\end{qoptions}
\qn{6} In the cubic structure of ZnS (Zinc blend) of edge length `$a$' :
\begin{qoptions}
\opt{A}{Distance between two nearest $\mathrm{Zn^{2+}}$ ions is $\dfrac a2$}
\opt{B}{Distance between two nearest $\mathrm{Zn^{2+}}$ ions is $\dfrac a{\sqrt2}$}
\opt{C}{The angle made by the lines connecting any $\mathrm{Zn^{2+}}$ to any two of its nearest neighbours is $60^\circ$}
\opt{D}{The angle made by the lines connecting any $\mathrm{Zn^{2+}}$ to any two of its nearest neighbours is $109^\circ28'$.}
\end{qoptions}
\qfullbegin
\sect{SECTION-I(ii)}
\paraq{7}{8}
Alcohols are converted to tosylates by treatment with p-toluene sulfonyl chloride (TsCl) in presence of pyridine. This overall process converts a poor leaving group ($\mathrm{O\overset{\ominus}{H}}$) into better one ($\mathrm{O\overset{\ominus}{T}s}$). A tosylate is a better leaving group because its conjugated acid p-toluene sulfonic acid is strong acid. Because alkyl tosylates have better leaving tendency they undergo both nucleophilic substitution and $\beta$-elimination.
\qfig{p1_chem_s2_para7_8_1}
\qfullsplit
\qfig{p1_chem_s2_para7_8_2}
$\mathrm{CH_3CH_2}$--$\mathrm{OTs}$ $\uparrow$ Better leaving group
\qfig{p1_chem_s2_para7_8_3}
\qfig{p1_chem_s2_para7_8_4}
\qfullend
\qn{7} Find the major product of following reaction :
\qfig{p1_chem_s2_q07_1}
\begin{qoptions}
\opt{A}{\qfigi{p1_chem_s2_q07_A}}
\opt{B}{\qfigi{p1_chem_s2_q07_B}}
\opt{C}{\qfigi{p1_chem_s2_q07_C}}
\opt{D}{\qfigi{p1_chem_s2_q07_D}}
\end{qoptions}
\qn{8}
\qfig{p1_chem_s2_q08_1}
Choose incorrect option for given reaction sequence
\begin{qoptions}
\opt{A}{Structure of product `$S$' have two stereoisomer}
\opt{B}{`$P$' and `$S$' are positional isomer of each other}
\opt{C}{`$P$' and `$S$' can be differentiated by Lucas test.}
\opt{D}{$P$, $Q$, $R$ and $S$ contains same number of chiral centre.}
\end{qoptions}
\paraq{9}{10}
Study the following thermochemical equations :\\
(a) $\mathrm{N_2O(g)+3H_2(g)\to H_2O(\ell)+N_2H_4(\ell)}$; $\Delta H=-76$ Kcal\\
(b) $\mathrm{4N_2(g)+3H_2O(\ell)\to 2NH_3(g)+3N_2O(g)}$; $\Delta H=+240$ Kcal\\
(c) $\mathrm{N_2H_4(\ell)+H_2O(\ell)\to 2NH_3(g)+\tfrac12O_2(g)}$; $\Delta H=+36$ Kcal\\
(d) $\mathrm{H_2O(\ell)\to H_2(g)+\tfrac12O_2(g)}$; $\Delta H=+68$ Kcal
\qn{9} The enthalpy of combustion (in Kcal/mol) of liquid $\mathrm{N_2H_4}$ is :
\begin{qoptions}
\opt{A}{$-129$}
\opt{B}{$-138$}
\opt{C}{$-147$}
\opt{D}{$+147$}
\end{qoptions}
\qn{10} The enthalpy of formation (in Kcal/mol) of liquid $\mathrm{N_2H_4}$ is :
\begin{qoptions}
\opt{A}{$-11$}
\opt{B}{$+44$}
\opt{C}{$+11$}
\opt{D}{$+13$}
\end{qoptions}
\columnbreak
\sect{SECTION-II}
\qn{1} In the acid catalysed hydrolysis of ethoxy ethanoate the rate law is found to be Rate $=K[\mathrm{CH_3COOC_2H_5}][\mathrm{H^+}]$. Calculate the value of ``$K$'' in $\mathrm{M^{-1}\,min^{-1}}$ from the following titration data.\par\smallskip
{\footnotesize\setlength{\tabcolsep}{3pt}\begin{tabular}{|p{3.0cm}|l|l|l|}
\hline
\multicolumn{4}{|l|}{$[\mathrm{H^+}]=0.1$ M}\\
\hline
Time (min) & $t=0$ & $t=6.93$ min & $t=\infty$\\
\hline
Volume of NaOH required for complete neutralization of small and equal volume of reaction mixture taken out at different time interval. & 35 ml & 60 ml & 85 ml\\
\hline
\end{tabular}}
\qn{2} If $X=$ Number of species among $\mathrm{O_2}$, $\mathrm{N_2}$, $\mathrm{B_2}$, $\mathrm{N_2^+}$, $\mathrm{O_2^{2-}}$, $\mathrm{C_2}$ in which HOMO-LUMO transition is either from BMO to BMO or ABMO to ABMO.\\
$Y=$ Number of planar species among $\mathrm{Be_2Cl_4}$, $\mathrm{I_2Cl_6}$, $\mathrm{Be(OH)_2(s)}$, $\mathrm{Al_2Cl_6}$, $\mathrm{Al_2(C_6H_5)_6}$, $\mathrm{B_2H_6}$\\
then $(X/Y)$ is \underline{\hspace{1cm}}.
\qn{3} The number of isomers of $\mathrm{[Pt(en)Cl_2(SCN)_2]}$ is\underline{\hspace{1cm}}.
\qn{4} Consider the following species.\\
Cresol, Anilene, Potassium ethyl Xanthate, $\mathrm{Na_2S}$, $\mathrm{CuSO_4}$, NaCN\\
$X=$ Number of species that can be used as activators in Froth floatation\\
$Y=$ Number of species that are used as stabilizers in froth floatation.\\
$Z=$ Number of species used as collectors in Froth floatation\\
then $\left(\dfrac{Y+Z}X\right)$ is\underline{\hspace{1cm}}.
\qn{5} Find total number of correct statements :\\
(a) R configuration of a compound may result in R and S configuration in an $\mathrm{S_N2}$ reaction.\\
(b) 1-Bromo-1-phenylethane shows greater reactivity towards $\mathrm{S_N1}$ and $\mathrm{S_N2}$ reaction compared to isopropyl bromide.\\
(c) Good yields of R--I may be obtained by heating alcohols with sodium or potassium iodide in 95\% orthophosphoric acid.\\
(d) All Lewis acid can act as electrophile and reverse is also true\\
(e) Nucleophilicity of $\mathrm{S\overset{\ominus}{H}}$ is greater than $\mathrm{O\overset{\ominus}{H}}$ in dimethyl formamide.\\
(f) Allyl bromide gives $\mathrm{S_N1}$ reaction but vinyl chloride does not.\\
(g) The solvolysis of 1-chloro-2-pentene is lesser than 1-chloropentane in water.\\
(h) $\mathrm{NaI+H_2SO_4}$ is not a good reagent to produce alkyl iodide from alcohol because $\mathrm{H_2SO_4}$ oxidises $\mathrm{I^\ominus}$ into $\mathrm{I_2}$.\\
(i) 1-methylcyclopentane form total ten monochlorined products after monochlorination.
\qn{6} (1) \qfig{p1_chem_s3_q06_1}
\qfig{p1_chem_s3_q06_2}
Number of sp carbon in product `$P$' is $q$ and number of delocalized $\pi$ electron pair in product `$Q$' is $p$ then find $p+q$
\qn{7} Number of reactions which are correctly matched with their major product?\\
(a) $\mathrm{R-Cl+NaI}\xrightarrow{\text{Acetone}}\mathrm{R-I}$\\
(b) $\mathrm{R-Cl+AgF}\xrightarrow{\text{DMF}}\mathrm{R-F}$\\
(c) $\mathrm{R-OH+NaBr}\xrightarrow{\mathrm{H_2O}}\mathrm{R-Br}$\\
(d) $\mathrm{R-OH+NaBr}\xrightarrow{\mathrm{H_3PO_4}}\mathrm{R-Br}$
\qfig{p1_chem_s3_q07_1}
\qn{8} The entropy change (in J/K) of 2.0 moles of an ideal gas whose adiabatic exponent $\gamma=1.50$, if a result of a certain process, the gas volume increased 2.0 times while the pressure dropped 4.0 times, is ($\ell$n $2=0.7$)
\partsec{3}{MATHEMATICS}
\sect{SECTION-I(i)}
\qn{1} Let $(1+\sqrt2)^n=A_n+B_n\sqrt2$, where $A_n$ and $B_n$ are rational numbers and $n\in N$, then the value of $\displaystyle\lim_{n\to\infty}\dfrac{A_n}{B_n}$ is equal to -\\
(where [.] denotes greatest integer function)
\begin{qoptions}
\opt{A}{$\left[\displaystyle\lim_{x\to0}\dfrac{\sin x}x\right]$}
\opt{B}{$\left[\displaystyle\lim_{x\to\infty}\dfrac{\sin x}x\right]$}
\opt{C}{$\displaystyle\lim_{x\to\frac\pi4}\dfrac{|\cos x-\sin x|}{|x-\frac\pi4|}$}
\opt{D}{$\left[\displaystyle\lim_{x\to0}\dfrac{\sin\sqrt2x}x\right]$}
\end{qoptions}
\qn{2} Given $\displaystyle\lim_{x\to\alpha}\dfrac{\sin([x^3]-[x]^3)}{(x-\alpha)^3}$ exists finitely, then the value of $\alpha$ can be (where $\alpha\in I$ \& [.] denotes greatest integer function)-
\begin{qoptions}
\opt{A}{0}
\opt{B}{1}
\opt{C}{$-1$}
\opt{D}{2}
\end{qoptions}
\qn{3} If the solution of the equation, $\sin(\tan^{-1}x)=\sqrt{4-\left(\sin(\cos^{-1}x)+\cos(\sin^{-1}x)\right)^2}$ is `$a$', then
\begin{qoptions}
\opt{A}{$\sin^{-1}a+\cos^{-1}a=\dfrac\pi2$}
\opt{B}{$2\sin^{-1}a+\cos^{-1}a=\dfrac\pi2$}
\opt{C}{$\sin^{-1}a+3\cos^{-1}a=\dfrac{3\pi}2$}
\opt{D}{$\tan^{-1}a+\cos^{-1}a=\dfrac\pi2$}
\end{qoptions}
\qn{4} A right angled triangle has one leg of length 1, another of length `$x$' \& hypotenuse of length `$y$'. The angle opposite to the side $x$ is $\theta$ then -
\begin{qoptions}
\opt{A}{$y^n-x^n=\dfrac{1-\sin^n\theta}{\cos^n\theta}$}
\opt{B}{$\displaystyle\lim_{\theta\to\frac\pi2}(y-x)=0$}
\opt{C}{$\displaystyle\lim_{\theta\to\frac\pi2}\left(y^2-x^2\right)=1$}
\opt{D}{$\displaystyle\lim_{\theta\to\frac\pi2}\left(y^3-x^3\right)$ does not exist.}
\end{qoptions}
\qn{5} If $f:A\longrightarrow A$, where $A$ is a set of non-negative integers and satisfying the condition.\\
(i) $x-f(x)=19\left[\dfrac x{19}\right]-90\left[\dfrac{f(x)}{90}\right]$, $\forall\,x\in A$\\
(ii) $1900<f(1990)<2000$\\
then $f(1990)$ equals. where [ ] denotes greatest integer function
\begin{qoptions}
\opt{A}{1994}
\opt{B}{1900}
\opt{C}{1904}
\opt{D}{1890}
\end{qoptions}
\qn{6} Which of the following is/are correct ($\log_{10}2=0.3010$, $\log_{10}3=0.4771$)
\begin{qoptions}
\opt{A}{Number of digits in $(243)^{10}$ is 24}
\opt{B}{Number of digits in $2^{100}$ is 31}
\opt{C}{If $\log_45=a$ and $\log_56=b$ then $\log_32$ is $\dfrac1{2ab-1}$}
\opt{D}{If $\log_45=a$ and $\log_56=b$ then $\log_32=2ab-1$}
\end{qoptions}
\sect{SECTION-I(ii)}
\paraq{7}{8}
Given the quadratic equation $x^2-(2K+2)x+K+3=0$ ($K\in R$) has roots $\alpha$ and $\beta$, then
\qn{7} If $\alpha<2<\beta$, then the least integral value of `$K$' is
\begin{qoptions}
\opt{A}{$-1$}
\opt{B}{$-2$}
\opt{C}{1}
\opt{D}{2}
\end{qoptions}
\qn{8} If $\alpha\ge2$, $\beta\ge2$ then the number of integral values of $K$ satisfying this condition is
\begin{qoptions}
\opt{A}{0}
\opt{B}{1}
\opt{C}{2}
\opt{D}{more than 2}
\end{qoptions}
\paraq{9}{10}
Let $f(x)$ and $g(x)$ be two quadratic polynomials $f(x)=ax^2+bx+c$, $g(x)=dx^2+ex+f$ where $a,b,c,d,e,f\in R$
\qn{9} If $a,b,c$ are the positive real numbers and $f(x)\ge0\ \forall\,x\in R$, then the minimum value of $\left(\dfrac{a+b+c}b\right)$ is
\begin{qoptions}
\opt{A}{1}
\opt{B}{2}
\opt{C}{3}
\opt{D}{0}
\end{qoptions}
\qn{10} If $H(x)=f(x).g(x)$ where $ad=1$ and $H(1)=1$, $H(2)=2$, $H(3)=3$, then $\left\{\dfrac{H(-1)+H(5)}{H(0)+H(4)}\right\}=$\\
(where $\{\ \}$ represents fractional function) is
\begin{qoptions}
\opt{A}{$\dfrac57$}
\opt{B}{$\dfrac27$}
\opt{C}{$\dfrac37$}
\opt{D}{$\dfrac47$}
\end{qoptions}
\sect{SECTION-II}
\qn{1} Let $f(x)=a_0+a_1x+a_2x^2+a_3x^3+a_4x^4$ where $a_i\in R$ and $\displaystyle\lim_{x\to0}\left(\dfrac{f(x)-x+2x^2}{3x^3}\right)^{\frac2x}=e^4$, then $\displaystyle\sum_{i=0}^{4}a_i$ is equal to
\qn{2} If $\alpha$ and $\beta$ are roots of equation $(\log_2x)^2-3\log_2x+4=0$, then the value of $\log_\alpha\beta+\log_\beta\alpha$ is equal to -
\qn{3} Consider a function\par\noindent\hfill$f(x)=\begin{cases}2x-1&-1\le x\le2\\x^2-3x+1&2<x\le3\end{cases}$\hfill\null\par
The Number of values of $x$ in the interval $[-3,3]$ satisfying $|f(|x|)|=1$, is
\qn{4} If $\displaystyle\lim_{x\to0}\dfrac{a+\log_e(1+x)+b\sin x}{x^2}=\ell$ (finite) then $|a+b+4\ell|$ (where $a,b,\ell\in R$)
\qn{5} Consider the equation $x^2+2x-n=0$, where $n\in N$ and $n\in[5,100]$. Total number of different values of `$n$' so that the given equation has integral roots, is
\qn{6} If a sequence is such that $a_1=1$ and $a_n=n(1+a_{n-1})\ \forall\,n\ge2$ and $n\in N$, then $\displaystyle\lim_{n\to\infty}\ell n\left(\left(1+\dfrac1{a_1}\right)\left(1+\dfrac1{a_2}\right)\cdots\left(1+\dfrac1{a_n}\right)\right)$ is \underline{\hspace{1.5cm}}.
\qn{7} If $\displaystyle\lim_{x\to0^+}\left(\dfrac{\sin^2x}{\left[\frac x\pi\right]+\frac{x^2}{\pi^2}}\right)\times\dfrac{\sin(\sin x)-\sin x}{ax^5+bx^3+c}=\dfrac{-\pi^2}{12}$. Then value of $b+c$ is\\
(Where [.] denotes greatest integer function)
\qn{8} If the value of $\displaystyle\lim_{x\to\infty}x^2\tan\left(\ell n\left(\sqrt{\cos\dfrac\pi x}\right)\right)$ is equal to $-k\pi^2$, then $k$ is equal to
\end{multicols}
\end{document}
